\chapter{32}



Von fern erklang das Schlagen von Kirchenglocken.

«Schon Mitternacht, kaum zu glauben, was für eine Nacht.»

Manuela sah nachdenklich aus. Sie richtete ihren Blick zum Himmel.

«Morgen ist Vollmond.»

Dabei streckte sie ihre Beine unter dem Tisch gerade aus und faltete die
Hände über ihrem Bauch. Joh lächelte sie an. Er begann die Schälchen der
Nachspeise zusammenzustellen und blieb, mit dem Geschirr beladen und dem
Blick zum Mond, einen runden Augenblick unbeweglich stehen.

Es fiel Louis zunehmend schwer, die Augen offen zu halten. Das Essen
hatte ihn ebenso ermüdet wie die Aufnahme der vielen Informationen zu
einer möglichen Anstellung. Ihre schnellen und präzisen Antworten auf
seine Zwischenfragen zeugten zweifellos von Erfahrung. Louis gefiel, was
er hörte. Den jungen Mann, David, den Louis hatte vorbeihuschen sehen,
erwähnte Joh nur kurz. Und schloss damit, dass er in Bezug auf seine
Klienten unter Schweigepflicht stehe.
Louis fragten sie wenig Persönliches. Ihr Interesse blieb immer mal
wieder bei seinem vermeintlichen Musikerleben hängen. Und am liebsten
wollten sie ihn spielen hören. Doch Louis winkte zwei Mal ab. Beim
dritten Mal spielte er den Verlegenen.

«Heute treff ich keine Töne mehr, aber morgen spiel ich euch gern was,
okay?»

«Ach ja, klar,» Joh legte seinen Kopf schräg, «doch bevor du vom Stuhl
fällst, Ciro,» seine Stimme klang augenblicklich sachlicher, «nach
unserem gemeinsamen Abend steht meiner Meinung nach weiteren Gesprächen
nichts im Weg,» jetzt blickte er fragend zu Manuela, die nickte, «wenn
du interessiert bist, sprechen wir morgen weiter?»

«Sehr gerne. Ich bin gespannt, worin meine konkrete Arbeit bei euch
aussehen könnte?»

Dann bedankte sich Louis für ihre freundliche Aufnahme und wünschte eine
gute Nacht. Er schleppte sich die Treppe hoch. Sympathische Leute,
dachte er. Wenn ihn sein Gefühl nicht täuschte, würde sein Plan
aufgehen.
Bevor er die Zimmertür schloss, hörte er Manuela und Johs Stimmen und
leises Geschirrklappern. Ein angenehmes Gefühl von Behaglichkeit umfing
ihn. Er atmete erleichtert auf. Wieder stieg ihm der Raumduft in die
Nase. Beim Anblick des Bettes, diesem frisch bezogenen, kleinen Wunder,
kamen ihm fast die Tränen. Und er dachte, dass sein Zustand auch etwas
mit seinem Alkoholpegel zu tun haben könnte, und dass er mehr als
erledigt war. Ausser seiner Hose, die er am Boden liegen liess, zog er
nichts aus. Er legte sich ins Bett und löschte das Licht.
Die Gespräche des Abends zogen noch einmal wie ein reich dekoriertes
Karussell an ihm vorbei. Mit der Aussicht auf einen zweiten Tag hier
oben, fühlte sich Louis erleichtert. Der würde mit Ausschlafen beginnen.
Joh hatte ihm dazu geraten.
Manuela wollte am nächsten Tag zusammen mit einer Freundin zum Gletscher
hochsteigen. Der wöchentliche Kontrollgang stehe an. Mazi, der Senn auf
der Alp oben, brauche frische Nahrungsmittel und Wäsche, hatten sie
erzählt. Auch die Salz-Lecksteine für die Schafe oben seien am Ausgehen.
Und Joh würde mit Louis einen möglichen Einsatz besprechen.
Schliesslich drehte sich Louis auf die Seite, zog die Decke bis dicht
ans Kinn hoch und fiel in den ersehnten Schlaf.

\sectionbreak

\indentedblock{«Es ist eiskalt. Festgepresster Schnee hat die Umgebung zugedeckt. Ein
Stück glatt abgeriegelte Welt liegt vor ihm. Weit, weit hinaus.

Da ist die junge Frau. Sie trägt ein rotes, kurzes Sommerkleid. Ist
barfuss.

Sie versucht auf einer Eisscholle zu stehen, die gefährlich hin und
herschwankt. Sie zittert. Das Meer ist mehrheitlich zugefroren. Die Frau
stellt sich direkt vor ihn, beinahe in Reichweite, ist nun eine
Pappfigur, flach, die nach hinten gebogen wird, um unverzüglich wieder
zurückzuschnellen, vor, zurück, vor, zurück, vor, zurück, im Takt.
Regelmässig.

Louis steht verzweifelt in einem hölzernen Ruderboot. Es schwankt
bedrohlich. Er sucht den Halt unter den Füssen. Und hört seine Stimme
ein verzweifeltes Nein schreien, im gleichen Rhythmus, immer wieder,
sobald sich die Frau nach hinten biegt und er seitlich in Schieflage
gerät. Nein, nein, bleib da, vor, zurück, vor zurück, vor. Er streckt
seine Hand nach ihr aus, verpasst jedes Mal knapp ihren Arm. Dann, dann
sieht er das Tier näherkommen, den geschmeidigen Tiger, sich leise
anschleichen. Der Tiger fletscht die Zähne und gibt ein Knurren von
sich. Louis versucht, ihn zu ignorieren. Der Tiger setzt zum Sprung an.
Er wird sich auf sie stürzen. Sie auffressen. Louis bebt vor Angst,
bettelt um Gnade.

Dann, der Schrei.»}

\sectionbreak

Louis erwachte. Nass vor Schweiss wähnte er sich in einem Feuerbecken 
gefangen. Er schlug sich die Bettdecke über den Kopf, um sie gleich
wieder hinunterzureissen, setzte sich auf und schnappte nach Luft. Er
hoffte, die Schreie nicht wirklich ausgestossen zu haben.
